\documentclass[conference]{IEEEtran}
\usepackage[utf8]{inputenc}
\usepackage[T1]{fontenc}
\usepackage{cite}
\usepackage{amsmath,amssymb}
\usepackage{graphicx}
\usepackage{booktabs}
\usepackage{url}

\title{Measuring the Marginal Detection Benefit of Sandboxed Emulation for LLM‑Generated NetOps Artifacts: A Paired Confirmatory Study}

\author{
\IEEEauthorblockN{Autonomous AI Researcher}
\IEEEauthorblockA{CNSM 2026 Agentic AI Researcher Track}
}

\begin{document}

\maketitle

\begin{abstract}
We preregistered and executed a paired confirmatory experiment (study\_id f2c3d8b1-7a6e-4a9a-b2c3-5f8e1d2b6c7e) that measured the marginal detection benefit of adding sandboxed emulation to a deterministic validator for LLM-generated intent\ensuremath{\rightarrow}configuration artifacts. Under shared\_initial\_candidate semantics using the declared model "gpt-5-mini" (provider: openai\_responses) we evaluated Npairs = 720 confirmatory tasks. The deterministic validator (deterministic\_netops\_validator\_v5) flagged 719/720 = 0.9986111111111111; the guarded pipeline (validator + sandboxed emulation) flagged 720/720 = 1.0. Paired difference (guarded - baseline) = 0.001388888888888884; exact McNemar two-sided p = 1.0; 95\% bootstrap percentile CI (10000 resamples, seed = 1729) = [0.0, 0.004166666666666667]. Run accounting: model\_calls\_used = 721 (720 shared initial + 1 repair-stage call). These artifact-grounded results lie far below our preregistered minimum effect-of-interest (0.20); we report the preregistered design, exact quantitative outcomes, run-accounting diagnostics, artifact-grounded interpretation for the negligible guarded benefit in this regime, and reproducibility guidance pointing to the archived execution bundle.
\end{abstract}

\section{Introduction}
Generative large language models (LLMs) are increasingly proposed to translate operator intent into device configuration and NetOps workflows. Operational deployment requires generated artifacts to satisfy syntactic correctness, workflow ordering, and higher-order safety/policy constraints. A common mitigation architecture is generate\ensuremath{\rightarrow}validate\ensuremath{\rightarrow}(repair)\ensuremath{\rightarrow}execute, sometimes augmented by sandboxed emulation to reveal runtime-only failures that static validators miss. Prior surveys and system reports advocate such workflow-centred mitigations [1]--[3], but provide limited large-scale paired measures of the marginal value added by sandboxed execution. Motivated by this gap, we preregistered and executed a controlled paired study to measure the aggregate detection-rate difference contributed uniquely by an automated sandboxed-emulation stage when the same initial generated candidate is analyzed by both conditions.

Research question and contributions. Our preregistered research question was: for synthetic intent\ensuremath{\rightarrow}configuration tasks under shared\_initial\_candidate semantics, what is the paired detection-rate difference (guarded - baseline) between (A) a deterministic validator-only pipeline and (B) the same deterministic validator followed by automated sandboxed emulation? Contributions: (1) a machine-readable preregistration and faithful autonomous execution; (2) an artifact-grounded paired analysis on Npairs = 720 with complete accounting; (3) an artifact-grounded interpretation explaining a near-zero marginal effect; (4) reproducibility guidance and prescriptive boundary conditions where guarded stages are likely to matter.

\section{Related Work}
Workflow-centred mitigations and verification tools. Recent survey and architectural reports document the practical trend of combining LLM generation with deterministic validators, emulators, and repair loops to raise operational assurance for NetOps artifacts [4]--[6]. These works emphasize that correctness in intent\ensuremath{\rightarrow}configuration pipelines is multidimensional: syntactic parsing, command-ordering/workflow constraints, invariants (e.g., keep‑one‑uplink‑up), and dynamic runtime properties may each require different V\&V primitives. Tool-focused case studies highlight the evolution and practical limits of static verifiers: mature configuration-analysis systems have progressively expanded rule-sets and topology-aware checks but still exhibit gaps for runtime-dependent behaviors and stateful sequences [7]. That evolution motivates layering an emulator after static validation: emulation can exercise temporal behaviors, device state transitions, and interactions that are difficult to capture exhaustively as static rules.

Empirical and prototype evidence. Prototype systems integrating generation with verification and execution-harnesses report heterogeneous gains: in some task regimes a verifier alone catches most dangerous errors, while in others emulation uncovers subtle runtime failures or environment-specific side effects that static rules miss [8], [9]. The practical implication drawn across this literature is that the marginal utility of a guarded (validator+emulation) pipeline depends on (a) validator coverage relative to the task distribution, (b) the fidelity of the emulator, and (c) the distributional properties of the input intents (e.g., constrained, parser-validated intents versus noisy natural-language requests). These prior observations motivated our paired experimental design: by reusing the identical initial generated candidate across conditions (shared\_initial\_candidate semantics) we isolate only downstream guarded-stage effects (emulation observations or permitted repair invocations) rather than conflating them with generation variability.

Gaps and evaluation needs. While surveys and system reports argue for generate\ensuremath{\rightarrow}validate\ensuremath{\rightarrow}repair and for emulator-assisted validation as best practice, they do not provide large-scale, preregistered paired estimates of the marginal detection benefit under controlled generation semantics. The literature also documents open questions our study probes: how much additional coverage does sandboxed emulation provide when a strong deterministic validator already targets the same workflow patterns, and which task regimes (adversarial inputs, cross‑vendor CLI heterogeneity, or stateful multi-step procedures) are most likely to reveal nontrivial guarded-stage benefits? By situating our paired confirmatory design amid these prior findings, we provide an artifact-grounded quantitative datapoint addressing one narrow slice of the broader evaluation agenda.

\section{Methodology}
Preregistration and confirmatory design. We preregistered a paired-binary confirmatory design (study\_id f2c3d8b1-7a6e-4a9a-b2c3-5f8e1d2b6c7e). The primary estimand is the paired\_success\_rate\_difference\_guarded\_minus\_baseline aggregated across confirmatory samples. The preregistered sample plan targeted Npairs = 720 (planned episodes = 1440) with a minimum effect-of-interest of 0.20 (20 percentage points) for power planning. The preregistration fixed inclusion/exclusion rules, missingness handling (deterministic reseed up to two times for parser failures; deterministic replacement for irrecoverable failures), multiplicity handling for exploratory secondary tests, and contamination controls based on deterministic hashing and version checks.

Task generation, topology templates, and vendor abstractions. Confirmatory items were deterministically generated from an intent\_configuration\_repair task family. The bank covers three abstracted vendor-CLI families and three topology templates (leaf--spine, 3-node service chain, triangle). Each task manifest entry encodes initial\_state, expected\_final\_state, explicit intent, required\_command\_sequence, allowed\_touched\_fields, and workflow policies (e.g., make-before-break). Inclusion required vendor-specific parser acceptance after up to two deterministic reseeds. In the confirmatory run there were zero irrecoverable parser exclusions and complete\_pair\_count = 720, consistent with the preregistered missingness policy.

Model, sampling semantics, and shared\_initial\_candidate enforcement. The declared/requested model was recorded as "gpt-5-mini" (provider: openai\_responses) and the model configuration archived (requested max\_output\_tokens = 2000; the client supplied no explicit temperature value). Important preregistered clarifications: (a) a null temperature field indicates the client supplied no explicit temperature; it does not imply provider-side determinism, and it does not alter the shared\_initial\_candidate contract; (b) shared\_initial\_candidate semantics were enforced per the preregistration: for each pair we obtained exactly one sampled initial candidate (a single provider call) and that identical candidate was presented to both downstream conditions (baseline and guarded). This design isolates downstream guarded-stage effects by removing generation variability as a source of between-condition discordance.

Baseline validator specification. The baseline validator (deterministic\_netops\_validator\_v5) implements deterministic, rule-based intent-constraint validation. Its core checks include: vendor-CLI syntactic parsing; per-command parsing and canonicalization; sequence-order checks against required\_command\_sequence; atomicity/workflow constraints (e.g., make-before-break); protected-interface invariants; group-level activity constraints; and final-state conformance tests. The validator returns a binary detection flag under a predeclared scoring rule and emits a validator\_trace object (validation\_before/after, parsed\_command\_count, required\_command\_count, violation\_codes).

Guarded pipeline and sandboxed-emulation harness. The guarded pipeline runs the baseline validator and then executes the candidate in an automated sandboxed-emulation harness that instantiates the declared topology template and executes normalized CLI commands against vendor-abstracted virtual devices. The emulation harness records runtime observations relevant to operational safety (interface admin transitions and timing, transient downtime and concurrent-interface states, MTU/VLAN application effects, and observable connectivity consequences). Emulation results are translated into runtime-observation rules that can trigger a guarded detection if they reveal policy or safety violations not encoded in the static validator. The emulation harness is conservative in mapping observations to detections; its purpose is to reveal runtime-only misconfigurations that the baseline validator might miss.

Repair policy and run accounting. The preregistered repair policy allowed at most one repair-stage model invocation per task; repair invocations were permitted only after the shared initial candidate was generated and after baseline validator and emulation observations were available for the guarded condition. Provider-call auditing for the confirmatory execution reports model\_calls\_used = 721 with stage\_counts showing shared\_initial\_generation = 720 and repair = 1; thus only a single repair-stage call occurred across the confirmatory set. Because shared\_initial\_candidate semantics were enforced, any discordant pair (guarded-only detection) is attributable to either an emulation-observed runtime violation or to the single documented repair invocation; both are recorded in the archived per-episode artifacts.

Inclusion, determinism, and contamination controls. The execution used deterministic artifact hashing, generator-seed recording, and explicit version checks before confirmatory runs. Parser-failure handling used deterministic reseeding with stored PRNG seeds under the task manifest; irrecoverable failures would be replaced deterministically and recorded. Deterministic reconciliation verifies complete\_pair\_count = 720, completed\_episode\_count = 1440, and failed\_episode\_count = 0. Provider-call auditing reports cache\_miss\_count = 721 and zero failed hosted model calls.

Analysis procedures and pre-registered inferential methods. The primary analysis is a paired-binary procedure aligned with McNemar-style inference. For the primary aggregated estimand we compute the paired detection-rate difference (guarded - baseline) with a two-sided exact McNemar test and produce a 95\% paired-bootstrap percentile confidence interval (bootstrap resamples = 10000, seed = 1729) as specified in the preregistration. Secondary per-class/vendor paired analyses were pre-specified and treated with Bonferroni-adjusted intervals. The primary analysis used the preregistered complete-pair handling; no irrecoverable pairs occurred, so conclusions rest on complete paired data.

Reproducibility guidance for independent verification. Per-episode raw scoring and trace artifacts were archived in the run's canonical execution bundle. Independent verifiers can locate discordant pairs by searching the run's archived raw results for entries where baseline\_score = 0 and guarded\_score = 1; the archived deterministic reconciliation and provider-call audit document the mapping from that pair to the single recorded repair-stage call. Summary paired counts and contingency outputs are archived with the analysis artifacts in the execution bundle. The archived artifacts were the direct input to the paired-binary executor and are the authoritative source for the numeric outputs reported below.

Design rationale and trade-offs. We chose shared\_initial\_candidate semantics to produce a conservative, direct estimate of guarded-stage marginal contributions: by reusing the same sampled initial candidate for both conditions, any observed discordance can only arise downstream of generation (emulation observation or repair). This isolates guarded-stage value but omits any potential benefit that could accrue from independent re-generation in the guarded condition (e.g., prompting-based repair prior to validator execution). Our inclusion policy (parser-validated synthetic tasks) reduces noise and ensures executability but concentrates the task bank where a strong validator can achieve high coverage; this choice is central to interpreting the near-zero aggregate effect observed in the confirmatory run.

\section{Results}
Primary confirmatory counts and test statistics (artifact-grounded). The analysis reports contingency counts: n11 = 719 (detected by both), n10 = 1 (guarded-only), n01 = 0 (baseline-only), n00 = 0 (neither). Marginal detection rates: baseline = 719/720 = 0.9986111111111111; guarded = 720/720 = 1.0. Paired estimate (guarded - baseline) = 0.001388888888888884. Exact McNemar two-sided test p = 1.0. Paired-bootstrap percentile 95\% CI (10000 resamples, bootstrap\_seed = 1729) = [0.0, 0.004166666666666667]. All numeric values below are reproduced verbatim from the run's analysis artifacts.

Table 1 --- Primary confirmatory summary.
- Npairs = 720 (complete paired episodes)
- Baseline successes = 719; Baseline success rate = 0.9986111111111111
- Guarded successes = 720; Guarded success rate = 1.0
- Paired difference (guarded - baseline) = 0.001388888888888884
- Exact McNemar two-sided p = 1.0
- 95\% bootstrap percentile CI = [0.0, 0.004166666666666667] (10000 resamples, seed = 1729)

Discordant-pair attribution and repair accounting. Provider-call accounting and deterministic reconciliation identify a single repair-stage invocation and the unique n10 discordant count. Because shared\_initial\_candidate semantics were enforced, the permissible causes for that single guarded-only detection are: (i) a guarded-stage emulation observation not encoded as a validator violation, or (ii) the single recorded repair-stage invocation. The archived raw-results record for the discordant entry contains the pair\_id and the validator\_trace and repair\_provider\_trace fields that link the guarded detection to the documented repair-stage call; these per-episode artifacts are included in the execution bundle for independent inspection.

Representative per-episode diagnostics. Representative scoring traces include validator\_trace objects with normalized\_configuration, validation\_before/after, parsed\_command\_count, required\_command\_count, and violation\_count. Those representative artifacts show that the vast majority of shared\_initial candidates complied with validator rules (violation\_count = 0) and that sandboxed emulation largely corroborated validator findings in this curated task bank.

Secondary analyses and missingness. The confirmatory run had complete\_pair\_count = 720 and no irrecoverable pairs. Secondary per-class/vendor analyses were pre-specified but are exploratory given the extreme marginal rates observed in the aggregated counts.

\section{Discussion}
Why was the guarded-stage aggregate benefit negligible? The archived evidence supports three artifact-grounded factors and suggests practical boundary conditions where guarded stages are---and are not---likely to matter.

1) Validator ceiling effect and task-rule alignment. deterministic\_netops\_validator\_v5 explicitly encodes the syntactic and semantic checks exercised by the curated tasks (representative patterns include shutdown\_configure\_restore, make\_before\_break\_uplink\_switchover, and paired\_link\_atomic\_mtu\_change). Validator traces show that most shared\_initial candidates satisfied the validator's rule set (violation\_count = 0). The baseline marginal success rate of 719/720 (0.9986111111111111) leaves essentially no headroom for the guarded emulation stage to flag additional failures in this constrained regime.

2) Curated inclusion and reduced adversarial exposure. The confirmatory sample was assembled from a deterministic generator and required vendor-parser acceptance before inclusion. This inclusion rule intentionally restricted confirmatory items to parser-validated, policy-constrained scenarios; as a result, many runtime-only faults that emulation detects in more diverse or adversarial settings are underrepresented.

3) Shared\_initial\_candidate semantics and limited repair activity. By design we reused the exact sampled initial candidate for both baseline and guarded conditions; this isolates downstream guarded-stage signals as the only allowable source of between-condition discordance. Provider-call auditing reports model\_calls\_used = 721 and stage\_counts: shared\_initial\_generation = 720, repair = 1. The single n10 discordant pair therefore corresponds to the single recorded repair-stage invocation or to a guarded-emulation observation not codified as a validator rule. The archived per-episode traces allow an auditor to locate that unique pair and inspect validator\_trace and repair\_provider\_trace to determine whether the repair altered the candidate or whether the emulation exposed a runtime-only violation.

Interpretation and operational implications. For environments that maintain a high-coverage static validator tailored to expected intent patterns, and where inputs are parser-validated and curated, adding an automated sandboxed-emulation stage may have negligible incremental detection benefit at scale. This does not imply sandboxing is unnecessary; rather, it implies a conditional efficiency boundary: when validator coverage closely matches expected task patterns and the input pipeline enforces strong prevalidation, guarded runtime checks will detect relatively few additional faults. Conversely, in operational regimes characterized by (a) noisy natural-language intents, (b) adversarially crafted misconfigurations, (c) broader uncurated vendor-CLI diversity, or (d) multi-step stateful interactions, the guarded stage can surface runtime semantics and side-effects that static validators cannot anticipate.

Practical recommendations grounded in the artifacts. The archived execution and analysis artifacts suggest three pragmatic paths depending on operator goals: (i) invest in richer static validators that capture more workflow-level constraints when low-latency, high-throughput intent\ensuremath{\rightarrow}config translation is the priority; (ii) deploy guarded sandboxing selectively for high-risk or high-impact changes where emulation fidelity matters; (iii) combine both approaches and use adversarial/fuzzing task banks to identify validator gaps---the archived traces provide concrete seed cases and validator-rule examples to guide such follow-up efforts.

Limitations rooted in the archived evidence. The near-zero aggregate effect we observed is specific to the preregistered task bank, the single validator implementation, the declared/requested model, and the shared\_initial\_candidate semantics. The run accounting (completed\_episode\_count = 1440; complete\_pair\_count = 720; model\_calls\_used = 721; repair = 1) permits independent verification that the observed discordance is localized to a single downstream event. Because the study was powered a priori for a minimum detectable paired difference of 0.20, the empirical observation (\ensuremath{\approx}0.0014) lies far below that sensitivity and must be interpreted as a narrowly scoped near-zero finding for this experimental regime---not as broad evidence that guarded stages are always unnecessary.

Where guarded stages are likely to add value. The diagnostics indicate three concrete extensions where guarded emulation should be retested: (1) adversarially generated tasks designed to exercise validator gaps (e.g., subtle state-dependent ordering errors); (2) expanded vendor-CLI families and syntactic/semantic diversity to challenge parser and validator coverage; (3) higher-fidelity emulation that models vendor-specific control-plane subtleties and long-lived state transitions. The archived per-episode traces and representative scoring artifacts provide seed cases and validator-rule examples to guide these follow-ups.

Concluding synthesis. The run's analysis artifacts permit a clear mapping from the reported aggregate counts to the single discordant episode and the recorded repair-stage call. These artifacts support a principled, evidence-led conclusion: in this preregistered, curated regime, a strong deterministic validator absorbed almost all detectable faults and the guarded sandboxed-emulation stage produced a vanishing marginal detection gain. Whether that conclusion generalizes to more adversarial or production-like regimes requires the targeted extensions described above.

\section{Limitations}
\begin{itemize}
\item Synthetic task bank and deterministic task generator produce limited ecological diversity and create a validator-ceiling effect; results may not generalize to noisier, adversarial, or production distributions.
\item Single declared/requested model family (gpt-5-mini) and single validator implementation (deterministic\_netops\_validator\_v5) limit generalizability across models and validator families.
\item Preregistered shared\_initial\_candidate semantics (one initial generation reused across paired conditions) isolate guarded-stage contribution but remove between-condition generation variability; this design choice constrains discordance sources to downstream guarded-stage observations or repair calls.
\item Sandboxed-emulation fidelity relative to vendor/OS production behaviors and large-scale stateful interactions is limited by the emulation harness used; higher-fidelity emulators could change outcomes in other regimes.
\item Study power planning targeted a minimum detectable paired difference of 0.20; the observed aggregate effect (\textasciitilde{}0.0014) lies far below that design sensitivity and should be interpreted as a narrowly scoped near-zero finding for this experimental regime rather than a general negative result for all NetOps workflows.
\item Provider-side nondeterminism and hidden caching remain residual uncertainties despite deterministic reconciliation procedures; these are documented in the archived reconciliation and provider-call audit artifacts.
\end{itemize}

\section{Conclusion}
We preregistered and executed a paired confirmatory study comparing a strong deterministic validator with the same validator augmented by sandboxed emulation on a curated synthetic intent\ensuremath{\rightarrow}configuration task bank. Over 720 paired tasks the guarded pipeline produced a negligible aggregate detection gain (0.001388888888888884) with exact McNemar p = 1.0 and 95\% bootstrap CI [0.0, 0.004166666666666667]. Run accounting shows one repair-stage invocation corresponding to the single guarded-only detection; because shared\_initial\_candidate semantics were enforced, archived per-episode traces permit independent verification of that mapping in the execution bundle. We interpret the result as arising from validator ceiling effects and curated task design; follow-up work should focus on adversarial and production-like task banks, multi-validator ensembles, and higher-fidelity emulation to identify regimes where guarded stages add substantive operational value.

Reproducibility pointers (compact). The canonical execution bundle archived by the run contains the execution manifest, the raw per-episode results file, the deterministic reconciliation / provider-call audit, the model configuration used for the run, and the analysis artifacts that include the paired contingency table and condition summary. These artifacts together reproduce the numeric outputs reported in this paper and allow an auditor to locate the unique discordant episode and to inspect validator and repair transcripts. The preregistration record is identified by study\_id f2c3d8b1-7a6e-4a9a-b2c3-5f8e1d2b6c7e and is recorded in the run's execution manifest for independent verification.

\section*{Disclosure Statement}
Disclosure: Autonomous post-lock research run executed under the frozen master prompt supplied to the system. Declared/requested model: "gpt-5-mini" (provider recorded as openai\_responses). Deterministic validator: deterministic\_netops\_validator\_v5. Orchestration adapter family: hosted\_netops\_gvr\_v1. Preregistration id: f2c3d8b1-7a6e-4a9a-b2c3-5f8e1d2b6c7e (preregistration record is archived in the run's execution manifest). Initial master-prompt immutable reference: provenance/master\_prompt.txt, SHA-256 = 1872df1e\allowbreak{}1805d2d9\allowbreak{}6940456c\allowbreak{}a016bd66\allowbreak{}5d1d5196\allowbreak{}add77f5a\allowbreak{}cdf1582b\allowbreak{}b39b15ba. Execution summary (artifact-grounded): completed\_episode\_count = 1440, complete\_pair\_count = 720, model\_calls\_used = 721 (720 shared initial + 1 repair-stage call). Human role and autonomy boundary: a human Research Director selected the broad NetOps topic family and constrained the pre-lock master prompt; after the autonomy lock the agentic pipeline autonomously performed literature selection, preregistration, code and prompt generation, experiment execution, analysis, peer review, and manuscript drafting. No manual human editing of technical manuscript content occurred after the autonomy lock. The run's execution manifest and archived artifacts serve as the canonical reproducibility bundle.

\begin{thebibliography}{99}
\bibitem{ref1} Muhammad Bilal, Jon Crowcroft, Ruizhi Wang, Xiaolong Xu, Schahram Dustdar, "Large Language Models for Agentic NetOps and AIOps: Architectures, Evaluation, and Safety", 2026, 10.48550/arxiv.2605.12729
\bibitem{ref2} Matt Brown, Ari Fogel, Daniel Halperin, Victor Heorhiadi, Ratul Mahajan, Todd Millstein, "Lessons from the evolution of the Batfish configuration analysis tool", 2023, 10.1145/3603269.3604866
\bibitem{ref3} http://arxiv.org/abs/2407.08249
\end{thebibliography}

\end{document}
